\chapter{A construção de uma transacção RingCT}
\label{chapter:transactions}

\notaFonte{%
  \fonteMargem{src/cryptonote_core/cryptonote_tx_utils.cpp}{cryptonote_tx_utils.cpp}%
  \codigoMargem{construct_tx_with_tx_key()}}%
Voltemos à saída de 11 unidades do capítulo~\ref{chapter:introduction}.
Alice conhece a chave que a pode gastar e a abertura do compromisso
que a acompanha. Quer que Bob receba 7, que lhe voltem 3 e que a taxa
seja 1. Já sabemos construir cada peça do pagamento. Falta descobrir
que dados passam de uma peça à seguinte para que o nó possa verificar
a transacção inteira.
\index{pagamento de Alice a Bob@pagamento de Alice a Bob}
\index{equilibrio de montantes@equilíbrio de montantes}
\index{prova de intervalo@prova de intervalo!num pagamento RingCT@num pagamento RingCT}

\section{As onze unidades vistas pelo nó}
\label{sec:pagamento-ringct}

Escrevamos \(C=xG+11H\) para o compromisso antigo. A carteira de
Alice usa o endereço de Bob para a primeira saída do exemplo e um
endereço seu para o troco. Os segredos partilhados de cada destino
dão as novas máscaras \(y_B,y_A\), segundo a
secção~\ref{sec:pedersen_monero}. O nó receberá os pontos
\[
 C_B=y_BG+7H,\qquad C_A=y_AG+3H.
\]
Os índices de transmissão podem surgir em qualquer das duas ordens.
Os índices \(A,B\) são aqui uma comodidade do leitor, que conhece os
destinos. Nenhum deles figura como nome no pagamento.

A máscara de \(C\) estava fixada muito antes de Alice pensar nesta
compra. Obtê-la somando \(y_B\) e \(y_A\) seria mera coincidência.
Alice dá, pois, à entrada outra apresentação, com
\(x'=y_B+y_A\pmod l\):
\[
                       \widetilde C=x'G+11H .
\]
\index{pseudo-compromisso@pseudo-compromisso!num pagamento RingCT@num pagamento RingCT}
O nó tem ambos os compromissos de saída, o pseudo-compromisso
\(\widetilde C\) e a taxa pública. Pode calcular \(C_B+C_A+H\)
e compará-lo com \(\widetilde C\). A razão por que os pontos
coincidem vê-se desenvolvendo a soma:
\begin{align*}
 C_B+C_A+H
  &=(y_B+y_A)G+(7+3+1)H\\
  &=x'G+11H=\widetilde C .
\end{align*}
O nó só teve de somar pontos e comparar o resultado. Os escalares \(x',y_B,y_A\) e os três montantes usados
para explicar a igualdade continuam nas carteiras.

Falta ligar o ponto do membro esquerdo ao dinheiro que Alice possuía.
O anel inclui \(C\) ao lado da chave da saída. Na posição que conhece,
Alice pode calcular a diferença de máscaras \(z=x-x'\). O ponto
correspondente é
\[
                 C-\widetilde C=(x-x')G+(11-11)H=zG.
\]
Assim se reencontram as duas testemunhas de CLSAG. Uma abre a chave
da saída, a outra dá o logaritmo desta diferença, ambas na mesma
coluna. Os pesos públicos associam as duas relações à lista concreta
dos candidatos e ao pseudo-compromisso. A chave de uso único dá
também a imagem que o nó procura no registo dos gastos. A
secção~\ref{sec:CLSAG} demonstrou por que o ciclo de desafios se
fecha com essas testemunhas.

Ainda poderia a igualdade dos compromissos aceitar, só pela soma,
quantidades impróprias para uma saída. No corpo \(\mathbb Z_l\),
\(12+(l-2)+1=l+11\) tem o mesmo resíduo que 11. O representante
\(l-2\) é, porém, muito maior que \(2^{64}-1\). A prova Bulletproofs+
exige uma decomposição em 64 bits para \emph{cada} uma das duas
saídas. Para os montantes honestos os primeiros bits, a começar nas
unidades, são \((1,1,1)\) e \((1,1)\), e todos os restantes são zero.
Dão \(1+2+4=7\) e \(1+2=3\). Um vector de 64 bits nunca recompõe
\(l-2\).

A prova de intervalo prende esses algarismos aos mesmos \(C_B,C_A\)
que entram na soma. A vinculação dos compromissos depende da
dificuldade de obter \(\log_G H\). No exemplo, mesmo contando duas
saídas e uma taxa, cada qual inferior a \(2^{64}\), o total é menor
que \(3\cdot2^{64}<l\). Entre somas inteiras nesse domínio, a
igualdade de resíduos implica igualdade de inteiros. O nó pôde assim
passar da igualdade de pontos à contabilidade do primeiro capítulo.
A assinatura, a prova de intervalo e o equilíbrio têm, cada qual,
um papel indispensável nessa passagem
\cite{MRL-0011-CLSAG,bulletproofs-plus-paper,monero-source-v16}.

\section{Formato RingCT com Bulletproofs+}
\label{sec:RCTTypeBulletproofPlus}

A versão de serialização é 2; o tipo RingCT é `6',
\texttt{RCTType}\allowbreak\texttt{BulletproofPlus}. Usa uma CLSAG por entrada, anéis de 16
membros e uma prova Bulletproofs+ agregada para as suas duas a dezasseis
saídas. A transacção de mineiro tem tipo `0', \texttt{RCTTypeNull}, e
montantes visíveis; a sua validade depende da entrada de geração, do
bloqueio temporal, dos alvos e da soma da recompensa. A construção de
mineiro considerada cria uma saída. Essa escolha de construção é distinta
das condições de validade, que verificam a estrutura e a soma dos montantes.

\emph{Confidencial} refere-se à ocultação dos montantes nos compromissos;
\emph{em anel} refere-se à autorização mediante uma posição secreta entre
saídas públicas. A taxa, as contagens de entradas e saídas, os membros dos
anéis, as imagens de chave e o campo \texttt{extra} permanecem públicos.

\subsection*{Formato histórico com Bulletproofs}
\label{sec:RCTTypeBulletproof2}

\notaFonte{%
  \fonteMargem{src/ringct/rctSigs.cpp}{rctSigs.cpp}%
  \codigoMargem{proveRctMGSimple()}\par \codigoMargem{verRctMGSimple()}\par \codigoMargem{verRctSemanticsSimple()}}%
O tipo `4', \texttt{RCTTypeBulletproof2}, combina Bulletproofs com MLSAG.
O apêndice \ref{appendix:RCTTypeBulletproof2} conserva uma transacção real
desse formato para a leitura dos bytes. O tipo `5', \texttt{RCTTypeCLSAG},
combina Bulletproofs com CLSAG. A leitura dos tipos históricos permite
verificar as transacções admitidas pelas regras das respectivas alturas;
na especialização da versão 16, as novas transacções ordinárias usam o
formato com Bulletproofs+ acima definido.

\subsection{Compromissos a montantes e taxas de transacção}
\label{sec:commitments-and-fees}
\index{equilibrio de montantes@equilíbrio de montantes}
\index{pseudo-compromisso@pseudo-compromisso!mascaras@máscaras}
\index{taxa de transaccao@taxa de transacção!equilibrio@equilíbrio}

Um remetente recebeu \(m\) saídas com montantes
\(a_1,\ldots,a_m\), endereços ocultos
\(P_{\pi_1,1},\ldots,P_{\pi_m,m}\) e compromissos
\begin{align*}
 C^a_{\pi_j,j}=x_jG+a_jH,\qquad j\in\{1,\ldots,m\}.
\end{align*}
Ele conhece as chaves privadas de gasto \(p_j\) tais que
\(P_{\pi_j,j}=p_jG\), os montantes \(a_j\) e as máscaras \(x_j\).

Pretende criar \(p\) saídas com montantes \(b_0,\ldots,b_{p-1}\) e pagar uma
taxa \(f\), todos em unidades atómicas, de modo que
\begin{align}
 \sum_{j=1}^{m}a_j
   =\sum_{t=0}^{p-1}b_t+f.
 \label{eq:transaction-amount-balance}
\end{align}
A taxa \(f\) é pública e serializada como um inteiro de comprimento
variável. O mínimo de taxa que os nós costumam exigir para propagar uma
transacção é política da carteira e da \emph{mempool}, não esta equação de
consenso: um bloco pode conter uma transacção de taxa inferior, desde que
esta seja válida nos restantes aspectos. O valor declarado continua, claro,
a entrar no equilíbrio e na recompensa do mineiro. Veja-se a secção
\ref{subsec:dynamic-minimum-fee}.

Para cada saída nova a carteira deriva, como na secção
\ref{sec:pedersen_monero}, uma máscara determinística
\begin{align*}
 y_t=\mathcal H_n(``\texttt{commitment\_mask}''\,\|\,h_t)
\end{align*}
e publica o compromisso
\begin{align*}
 C^b_t=y_tG+b_tH.
\end{align*}

Para cada entrada cria também um pseudo-compromisso
\begin{align*}
 \widetilde C^a_j=x'_jG+a_jH.
\end{align*}
As primeiras \(m-1\) pseudo-máscaras \(x'_j\) são escolhidas ao acaso. A
última é determinada pela condição de balanço:
\begin{align*}
 x'_m=\sum_{t=0}^{p-1}y_t-\sum_{j=1}^{m-1}x'_j\pmod l.
\end{align*}
Por homomorfismo dos compromissos, a equação
\begin{align}
 \sum_{j=1}^{m}\widetilde C^a_j
 \stackrel{?}{=}
 \sum_{t=0}^{p-1}C^b_t+fH
 \label{eq:transaction-commitment-balance}
\index{c tilde@\(\protect\widetilde C\), pseudo-compromisso}
\end{align}
verifica ao mesmo tempo o equilíbrio dos montantes e o das máscaras, sem
revelar nenhum \(a_j\) ou \(b_t\). A própria igualdade não prova que o autor
possui as entradas. É a CLSAG que prende cada pseudo-compromisso à saída
real de um anel.

O programa de referência baralha as saídas antes de lhes atribuir os índices
\(t\); este baralhamento é uma escolha da carteira. As entradas, pelo contrário,
têm de estar ordenadas por imagem de chave em ordem estritamente decrescente
dos bytes da codificação. Desde a versão 7, o consenso rejeita uma transacção
que não respeite essa ordem.


\subsection{Dados de cada saída}

Para a saída \(t\), o remetente escolhe a chave de transacção apropriada e
calcula a derivação partilhada \(D_t\), o escalar
\begin{align*}
 h_t=\mathcal H_n(D_t\,\|\,\operatorname{varint}(t))
\end{align*}
e o endereço oculto do destinatário
\begin{align*}
 P^b_t=h_tG+K^s_t.
\end{align*}
Na versão 16 o alvo serializado é etiquetado: contém \(P^b_t\) e o primeiro
byte de
\begin{align*}
 \mathcal H(``\texttt{view\_tag}''\,\|\,D_t\,\|\,
                    \operatorname{varint}(t)).
\end{align*}
A etiqueta deixa a carteira rejeitar quase todas as saídas alheias antes de
fazer o trabalho mais caro; não é uma prova de propriedade.

\notaFonte{%
  \fonteMargem{src/device/device_default.cpp}{device_default.cpp}%
  \fonteMargem{src/ringct/rctOps.cpp}{rctOps.cpp}%
  \fonteMargem{src/ringct/rctTypes.h}{rctTypes.h}}%
O montante cifrado contém apenas oito bytes:
\begin{align*}
 \operatorname{LE8}(b_t)\mathbin\oplus
 \mathcal H(``\texttt{amount}''\,\|\,h_t)[0\ldots7].
\end{align*}
A máscara não é transmitida cifrada ao lado do montante. O destinatário volta a
calcular \(y_t\) pela fórmula anterior e confirma que o compromisso publicado
é \(C^b_t=y_tG+b_tH\). Esta confirmação é importante:
decifrar oito bytes com aspecto razoável não basta para possuir uma saída.


O campo \texttt{extra} contém a chave pública principal da transacção.
Para os casos com sub-endereços, pode conter também uma lista de
chaves adicionais, uma por saída, conforme os casos da secção
\ref{sec:multi_out_transactions}. Essas chaves permitem ao destinatário
obter \(D_t\); não demonstram ao consenso que certo endereço foi o
destinatário pretendido.


\subsection{Uma prova Bulletproofs+ para todas as saídas}

Os compromissos \(C^b_0,\ldots,C^b_{p-1}\) entram numa única prova
Bulletproofs+ que mostra
\begin{align*}
 b_t\in[0,2^{64}-1]\qquad\text{para }0\leq t<p.
\end{align*}
Na transacção ordinária v16, \(2\leq p\leq16\). A prova é preenchida
internamente até
\begin{align*}
 M=2^{\lceil\log_2p\rceil},
\end{align*}
e o verificador exige precisamente o menor \(M\) que cobre as saídas; não se
pode escolher uma prova maior só para mudar o aspecto dos bytes.

\notaFonte{%
  \fonteMargem{src/ringct/rctSigs.cpp}{rctSigs.cpp}%
  \fonteMargem{src/ringct/rctTypes.cpp}{rctTypes.cpp}%
  \fonteMargem{src/cryptonote_basic/cryptonote_format_utils.cpp}{cryptonote_format_utils.cpp}%
  \fonteMargem{src/ringct/bulletproofs_plus.cc}{bulletproofs_plus.cc}}%
O objecto serializado é
\begin{align*}
 \Pi_{\mathrm{BPP+}}=
 (A,A_1,B,r_1,s_1,d_1,\mathbf L,\mathbf R).
\end{align*}
Os compromissos \(V_t\) não são serializados dentro da prova: o nó
reconstrói-os a partir dos \(C^b_t\) que já estão na base RingCT. Aplica-se a
seguinte convenção de cofactor. \texttt{outPk} guarda o
compromisso normal \(C^b_t\), enquanto a prova usa
\(V_t=8^{-1}C^b_t\). Do mesmo modo, \(A,A_1,B,L_i,R_i\) são serializados
divididos por 8 e o verificador multiplica-os por 8. Esta codificação não
altera o montante comprometido. A prova e a verificação completas estão na
secção \ref{sec:bulletproof_plus}. Os compromissos históricos do anel, os
compromissos de saída em \texttt{outPk} e os pseudo-compromissos são
serializados na representação normal; não se lhes aplica esta divisão por 8.



\subsection{Um anel: uma entrada real e quinze desvios}

Para cada entrada \(j\), a transacção contém 16 índices de saídas históricas.
Depois de converter os \emph{offsets} relativos em índices absolutos, o nó
obtém o anel público
\begin{align}
 \mathcal R_j=
 \bigl((P_{0,j},C^a_{0,j}),\ldots,(P_{15,j},C^a_{15,j})\bigr).
 \label{eq:transaction-ring}
\index{pi@\(\pi_j\), posição secreta no anel}
\index{anel@anel!lista ordenada@lista ordenada}
\index{anel@anel!posicao secreta@posição secreta}
\end{align}
Cada membro tem a mesma representação pública: uma chave de utilização
única \(P_{i,j}\) e o compromisso \(C^a_{i,j}\) da saída correspondente.
O índice \(i\) indica o lugar no anel ordenado. O remetente conhece a
abertura e a chave privada de um desses pares, no lugar secreto
\(\pi_j\); os outros quinze servem de desvios. Assim, \(P_{\pi_j,j}\)
é a própria chave pública que aparece na posição \(\pi_j\) da lista.
A informação adicional do dono está nos seus segredos.

O segundo índice, \(j\), identifica a entrada. Duas entradas têm anéis
próprios e podem ter posições secretas distintas, por exemplo
\(\pi_1=3\) e \(\pi_2=11\). A CLSAG de cada entrada relaciona a chave
de gasto e a diferença de máscaras na sua posição \(\pi_j\). Esta
coincidência é exigida dentro de cada anel; a transacção reúne as
autorizações de todos eles.

Na versão 16 todas as entradas RingCT têm exactamente quinze desvios, isto
é, dezasseis membros, e todas têm o mesmo tamanho de anel. Isto é consenso.
\notaFonte{%
  \fonteMargem{src/wallet/wallet2.cpp}{wallet2.cpp}%
  \codigoMargem{gamma_picker::pick()}\par \codigoMargem{get_outs()}}%
A escolha de \emph{quais} saídas ocupam os quinze lugares é comportamento da
carteira. O programa de referência insere a saída real e amostra candidatos
históricos com a sua \texttt{gamma\_picker}: parte de uma distribuição gamma
de forma \(19{,}28\) e escala \(1/1{,}61\), corrige a janela de maturação,
estima um índice pela densidade recente de saídas e escolhe ao acaso uma
saída dentro do bloco encontrado.\footnote{A distribuição dos desvios é uma heurística de privacidade da
carteira. A validade depende das referências e das provas. As imagens de
chave registam os gastos; a lista de blocos conserva as saídas RingCT sem
lhes atribuir uma marca pública de «gasta».}

Quem verifica conhece todos os pares de \(\mathcal R_j\), mas não recebe
esses pontos repetidos dentro da assinatura. Recebe os \emph{offsets} e vai
buscá-los à lista de blocos. Antes de verificar a CLSAG confirma que as referências
existem, têm tipos admissíveis, já maturaram e estão desbloqueadas.

\subsubsection{Dos índices globais aos deslocamentos relativos}
\label{subsec:global-output-indices}
\index{indice@índice!global de saida@global de saída}
\index{indice@índice!local de saida@local de saída}
\index{indice@índice!relativo@relativo}
\index{deslocamento relativo@deslocamento relativo}
\index{key offsets@\texttt{key\_offsets}}
\index{vin@\texttt{vin}}

Fixemos um prefixo da lista de blocos já aceite. As saídas são registadas
pela ordem dos blocos, das transacções em cada bloco e das saídas em cada
transacção. A numeração começa em zero. Nas referências RingCT, o campo
de montante visível da entrada é zero e selecciona a lista comum das
saídas registadas nessa classe. Cada saída recebe aí um \emph{índice
global}, também chamado índice absoluto. A numeração pertence à lista
inteira, atravessando as fronteiras entre transacções.

Esta lista inclui as saídas de mineiro de transacções de versão 2,
registadas com o compromisso \(G+bH\), de montante público \(b\) e
máscara igual a um. As antigas classes de montante visível
têm numerações próprias. Por isso, a identificação completa de uma
referência é o par formado pela classe de montante e pelo índice; nas
entradas RingCT aqui estudadas, a primeira componente vale sempre zero.

Há também o índice local \(t\) de uma saída dentro da transacção que a
criou. É esse \(t\), já usado na derivação do endereço oculto, que
acompanha o segredo partilhado no cálculo de \(h_t\). O índice global
é atribuído quando a saída entra na lista de blocos. Por exemplo, uma
saída com índice local \(t=1\) pode receber o índice global \(150\) e
ocupar depois o lugar \(i=2\) num anel. Os três números descrevem três
posições do mesmo objecto. Um nó pode recuperar a transacção de origem e
o índice local a partir do registo global. A correspondência é relativa
ao prefixo fixado; uma reorganização exige reconstruir a parte afectada.

Para uma entrada \(j\), sejam
\[
 0\leq g_{0,j}<g_{1,j}<\cdots<g_{15,j}
\]
os índices globais escolhidos. A sua codificação por deslocamentos é
\begin{equation}
 d_{0,j}=g_{0,j},\qquad
 d_{i,j}=g_{i,j}-g_{i-1,j}\quad(1\leq i<16).
 \label{eq:relative-output-indices}
\end{equation}
O vector \(\texttt{key\_offsets}=(d_{0,j},\ldots,d_{15,j})\) pertence
à entrada \(j\) de \texttt{vin}. O primeiro termo mede a distância à
origem; cada termo seguinte mede a distância ao membro anterior.
A operação inversa é a soma acumulada
\begin{equation}
 g_{i,j}=\sum_{r=0}^{i}d_{r,j}.
 \label{eq:absolute-output-indices}
\end{equation}
De facto, as diferenças intermédias cancelam-se:
\(g_{0,j}+(g_{1,j}-g_{0,j})+\cdots+(g_{i,j}-g_{i-1,j})=g_{i,j}\).
As contas são feitas com índices inteiros não negativos, representáveis
em 64 bits. Para a lista estritamente crescente,
\(d_{i,j}>0\) quando \(i>0\); o primeiro deslocamento pode ser zero,
referindo a primeira saída da classe. O teste de repetição aplica-se
aos deslocamentos posteriores ao primeiro.

Consideremos quatro posições para acompanhar a conta, conservando em
mente que o anel completo desta construção tem dezasseis membros:
\[
 (7,130,150,400)\ \longmapsto\ (7,123,20,250).
\]
A leitura inversa dá, sucessivamente, \(7\), \(7+123=130\),
\(130+20=150\) e \(150+250=400\). Cada entrada começa uma nova soma
acumulada. A ordenação das próprias entradas pelas imagens de chave,
já estabelecida acima, conserva dentro de cada uma o anel e a sua
assinatura associados.

\subsubsection{Dos bytes aos membros do anel}
\label{subsec:ring-reconstruction}
\index{anel@anel!reconstrucao@reconstrução}
\index{varint@varint}
\index{inteiro de comprimento variavel@inteiro de comprimento variável|see{varint}}

\notaFonte{%
  \fonteMargem{src/common/varint.h}{varint.h}}%
Cada deslocamento é escrito como um inteiro de comprimento variável,
ou \emph{varint}. Para \(n>0\), escreva-se
\[
 n=a_0+2^7a_1+\cdots+2^{7k}a_k,
 \qquad 0\leq a_r<128,\quad a_k>0.
\]
Os bytes surgem pela ordem \(a_0,\ldots,a_k\). Nos primeiros \(k\),
o bit superior vale um para indicar continuação; no último vale zero.
O inteiro zero escreve-se com um único byte zero. Esta é a representação
canónica: o último grupo de uma codificação com vários bytes contém
um algarismo não nulo. O comprimento é um para zero e
\(1+\lfloor\log_{128}n\rfloor\) para \(n>0\).

Assim, \(250=122+128\cdot1\). O primeiro byte contém
\(122+128=250=\texttt{FA}_{16}\), e o segundo contém
\(1=\texttt{01}_{16}\). No exemplo anterior, o vector completo,
incluindo o seu comprimento, é
\[
 \underbrace{\texttt{04}}_{\text{quatro termos}}\quad
 \underbrace{\texttt{07}}_{7}\quad
 \underbrace{\texttt{7B}}_{123}\quad
 \underbrace{\texttt{14}}_{20}\quad
 \underbrace{\texttt{FA 01}}_{250}.
\]
São seis bytes. Codificar directamente os quatro índices absolutos
exigiria oito bytes, incluindo o mesmo comprimento. A economia resulta
de muitos intervalos entre saídas caberem em menos grupos de sete bits
do que os índices acumulados. Para dezasseis membros, o prefixo do
vector é \texttt{10}, a representação hexadecimal do inteiro dezasseis.

Dentro de \texttt{vin}, uma entrada desta espécie segue a disposição
\[
 \underbrace{\texttt{02}}_{\text{tipo de entrada}}\ \|
 \underbrace{\texttt{00}}_{\text{classe de montante}}\ \|
 \operatorname{varint}(16)\ \|
 \operatorname{varint}(d_{0,j})\ \|\ \cdots\ \|
 \operatorname{varint}(d_{15,j})\ \|\ \operatorname{enc}(I_j).
\]
Os 32 bytes finais são a imagem de chave. A marca \texttt{02} identifica
o tipo de entrada; a versão da transacção é outro campo do prefixo.
O pseudo-compromisso e a CLSAG aparecem mais adiante, na parte RingCT.

O nó lê os dezasseis deslocamentos, calcula os índices absolutos e,
para cada \(g_{i,j}\), consulta o registo da classe zero. Este fornece
o par \((P_{i,j},C^a_{i,j})\), acompanhado dos dados necessários à
verificação do desbloqueio e da maturidade. A consulta conserva a
ordem dos índices; obtém-se assim exactamente a lista ordenada
\(\mathcal R_j\) da Equação~\eqref{eq:transaction-ring}.
Subtrair \(\widetilde C^a_j\) a cada compromisso prepara então a
segunda fila da CLSAG.

Toda esta reconstrução usa dados públicos. A posição secreta
\(\pi_j\) pertence à testemunha da assinatura: o verificador dispõe
do mesmo tipo de par em todas as posições. Os índices e os pontos
recuperados determinam a declaração pública à qual a CLSAG fica
vinculada. A imagem \(I_j\) fornece, separadamente, o marcador de
gastos repetidos. Uma saída pode figurar como desvio em vários anéis;
esse reaparecimento é compatível com a função pública da lista de
candidatos.


\subsection{Assinatura}
\label{full-signature}
\index{clsag@CLSAG!autorizacao de uma entrada@autorização de uma entrada}

O pseudo-compromisso da entrada \(j\) é subtraído a cada compromisso do
anel:
\begin{align*}
 Q_{i,j}=C^a_{i,j}-\widetilde C^a_j.
\end{align*}
Para os desvios \(i\ne\pi_j\), \(Q_{i,j}\) é apenas um ponto para o qual o
signatário não conhece necessariamente logaritmo em base \(G\). Na coluna
real, os montantes cancelam:
\begin{align}
 Q_{\pi_j,j}
 &=(x_jG+a_jH)-(x'_jG+a_jH)\notag\\
 &=(x_j-x'_j)G=z_jG.
 \label{eq:transaction-zero-commitment}
\end{align}

Assim, para assinar a entrada \(j\), o remetente conhece dois testemunhos na
mesma coluna secreta:
\begin{align*}
 p_j&=\log_G(P_{\pi_j,j}),\\
 z_j&=\log_G(Q_{\pi_j,j})=x_j-x'_j\pmod l.
\end{align*}
A CLSAG comprime estas duas relações numa só assinatura, sem revelar
\(\pi_j\). Com \(J_{i,j}=\mathcal H_p^{(1)}(P_{i,j})\), as duas
imagens usadas na agregação são
\begin{align*}
 I_j&=p_jJ_{\pi_j,j},\\
 D_j&=z_jJ_{\pi_j,j}.
\end{align*}
\(I_j\) é guardada na entrada, antes da parte RingCT. A assinatura CLSAG
serializa, para um anel v16, dezasseis respostas
\(s_{0,j},\ldots,s_{15,j}\), o desafio inicial \(c_{1,j}\) e
\begin{align*}
 \bar D_j=8^{-1}D_j.
\end{align*}
Aqui a primeira coluna tem índice zero. O nó começa o ciclo da
secção~\ref{sec:CLSAG} com \(c_{1,j}\) e \(s_{0,j}\). Depois da ronda
de \(s_{15,j}\), deve reencontrar \(c_{1,j}\).

Em suma, a \emph{declaração pública} de uma CLSAG desta transacção contém
\(\mathfrak m\), os dezasseis pares
\((P_{i,j},C^a_{i,j})\), \(\widetilde C^a_j\), \(I_j\) e \(\bar D_j\).
Os \emph{testemunhos privados} são o índice \(\pi_j\), a chave de utilização
única \(p_j\) e a diferença de máscaras \(z_j=x_j-x'_j\). Os \emph{hashes} de
agregação incluem ambas as filas, as imagens e o pseudo-compromisso; o ciclo
de desafios usa os dois testemunhos na mesma coluna. Esta coincidência de
índice vincula a chave de gasto à diferença correcta de compromissos.

Para \(m\) entradas, a assinatura total é
\begin{align*}
 \tau=\{\sigma_1(\mathfrak m),\ldots,\sigma_m(\mathfrak m)\}.
\end{align*}
Cada entrada tem o seu próprio anel e o seu próprio índice real. Todas as
CLSAG assinam a mesma mensagem \(\mathfrak m\), mas cada uma prende também o
seu pseudo-compromisso à sua assinatura. A transacção só passa se todas
passarem.


\subsection{A mensagem que as CLSAG assinam}

\notaFonte{%
  \fonteMargem{src/ringct/rctSigs.cpp}{rctSigs.cpp}%
  \codigoMargem{get_pre_mlsag_hash()}\par \codigoMargem{genRctSimple()}\par \codigoMargem{CLSAG_Gen()}%
  \fonteMargem{src/device/device_default.cpp}{device_default.cpp}}%
A mensagem comum das CLSAG reúne três resumos dos dados da transacção:
\begin{align}
 \mathfrak m=\mathcal H(h_0\,\|\,h_1\,\|\,h_2),
 \label{eq:transaction-pre-clsag-hash}
\index{transaccao@transacção!mensagem assinada@mensagem assinada}
\end{align}
onde:

\begin{description}
    \item[\(h_0\)] é o \emph{hash} do prefixo serializado: versão,
    \texttt{unlock\_time}, entradas (montantes visíveis zero, \emph{offsets}
    e imagens de chave), saídas (montantes visíveis zero, tipo de alvo,
    endereços ocultos e etiquetas de vista) e os bytes exactos de
    \texttt{extra};

    \item[\(h_1\)] é o \emph{hash} da base RingCT serializada: tipo `6',
    taxa, os oito bytes cifrados de cada montante e os compromissos normais
    de saída em \texttt{outPk};

    \item[\(h_2\)] é o \emph{hash} da prova Bulletproofs+ pela ordem
    \(A,A_1,B,r_1,s_1,d_1,L_0,\ldots,L_q,R_0,\ldots,R_q\).
\end{description}

Os corpos CLSAG ficam necessariamente de fora, pois ainda estão a ser
calculados. Há uma subtileza menos óbvia: os pseudo-compromissos também não
entram em \(h_0,h_1,h_2\). Em vez disso, a transcrição de cada CLSAG inclui
explicitamente o seu próprio \(\widetilde C^a_j\), juntamente com
\(\mathfrak m\), os \(P_{i,j}\), os \(C^a_{i,j}\), \(I_j\) e \(\bar D_j\). Logo,
mudar um pseudo-compromisso invalida a assinatura correspondente; não é, porém,
verdade literal que exista um único \emph{hash} de todos os bytes excepto as
assinaturas.

Esta disposição prende os destinatários, os anéis, as imagens de chave, a
taxa, \texttt{extra}, os montantes em código, os compromissos de saída, a
prova de domínio e cada pseudo-compromisso à autorização do dono. Uma
alteração deixa de fechar pelo menos um ciclo CLSAG.

Ao receber o pagamento, o nó verifica o intervalo de cada montante novo
com Bulletproofs+, a autorização de cada entrada com a respectiva CLSAG
e a conservação da soma com a equação de equilíbrio. Consulta ainda as
imagens de chave para detectar um gasto repetido. A validade exige estes
testes, a serialização correcta e as demais condições de consenso.


\subsection{Verificação}
\label{full-verify}

A implementação divide a verificação RingCT em duas partes com nomes um
pouco enganadores:

\begin{enumerate}
    \item a verificação \emph{semântica} confirma a forma dos vectores, a
    correspondência entre saídas e montantes em código, a existência de uma
    única prova Bulletproofs+ com a capacidade exacta, a equação
    \eqref{eq:transaction-commitment-balance} e a prova de domínio;

    \item a verificação \emph{não semântica} reconstrói os anéis a partir da
    lista de blocos e verifica uma CLSAG por entrada contra o seu
    pseudo-compromisso.
\end{enumerate}

\notaFonte{%
  \fonteMargem{tests/core_tests/bulletproof_plus.cpp}{bulletproof_plus.cpp}%
  \fonteMargem{tests/core_tests/rct2.cpp}{rct2.cpp}%
  \fonteMargem{tests/unit_tests/bulletproofs_plus.cpp}{bulletproofs_plus.cpp}%
  \fonteMargem{tests/unit_tests/serialization.cpp}{serialization.cpp}}%
Em redor destas duas partes, o núcleo verifica ainda a versão e o tipo, os
alvos das saídas, o número e a igualdade dos tamanhos dos anéis, os
\emph{offsets}, a existência e maturidade das referências, as imagens de
chave e o limite de peso. Das duas verificações, só a segunda exige a
consulta das saídas históricas para reconstituir os anéis.

Os escalares de resposta e desafio CLSAG têm de estar em forma canónica. Os
pontos têm de descodificar, \(I_j\) não pode ser a identidade e tem de
satisfazer \(lI_j=0\); depois de recuperar \(D_j=8\bar D_j\), a verificação
rejeita também a identidade. Bulletproofs+ verifica os escalares canónicos,
a forma dos vectores e as suas próprias equações depois das multiplicações
por 8. As condições sobre os pontos combinam, portanto, testes de domínio
e multiplicações pelo cofactor, com regras próprias em cada fronteira.

\subsection{Evitar o gasto duplo}

\notaFonte{%
  \fonteMargem{src/cryptonote_core/blockchain.cpp}{blockchain.cpp}%
  \codigoMargem{have_tx_keyimges_as_spent()}}%
A imagem de chave \(I_j=p_j\mathcal H_p(P_{\pi_j,j})\) é determinística para
a chave privada de utilização única, mas não revela qual dos dezasseis
\(P_{i,j}\) lhe corresponde. O núcleo rejeita imagens repetidas dentro da
mesma transacção, entre transacções do mesmo bloco e contra o conjunto de
imagens já gastas na lista de blocos. Sob a segurança da CLSAG e da função de
\emph{hash} para ponto, uma segunda assinatura com a mesma chave produz a
mesma \(I_j\): o marcador denuncia a repetição sem desmascarar o membro real.

A igualdade de uma imagem de chave não diz quanto foi gasto, para onde foi,
nem qual membro do anel era verdadeiro. Diz apenas que duas tentativas não
podem ambas entrar na história canónica. Antes disso, a \emph{mempool}
também evita aceitar conflitos que conhece; a regra decisiva é a validação
do bloco.


\subsection{Maturação e tempo de desbloqueio}

Todas as saídas referidas por um anel v16, verdadeira e desvios, têm de ter
pelo menos dez blocos. Uma saída criada no bloco de altura \(h\) pode, por
esta regra, aparecer pela primeira vez num anel de uma transacção incluída à
altura \(h+10\). Não é apenas o comportamento por omissão da carteira: desde
v12 o núcleo impõe-o ao conjunto inteiro de referências.

\notaFonte{%
  \fonteMargem{src/cryptonote_core/blockchain.cpp}{blockchain.cpp}%
  \codigoMargem{is_tx_spendtime_unlocked()}%
  \fonteMargem{src/cryptonote_core/tx_pool.cpp}{tx_pool.cpp}}%
O prefixo contém ainda um único \texttt{unlock\_time} para todas as saídas
da transacção. Valores inferiores a \(500\,000\,000\) são interpretados como
altura; valores a partir daí, como tempo Unix. Na versão 16, o teste temporal
usa o tempo ajustado derivado dos blocos anteriores e uma tolerância de 120
segundos. A carteira de referência constrói transacções ordinárias com
\texttt{unlock\_time}=0, e os nós não propagam normalmente uma transacção
ordinária com outro valor; isto é política de \emph{mempool}. O consenso
conserva a funcionalidade e verifica também o \texttt{unlock\_time} original
de cada membro do anel.

\notaFonte{%
  \fonteMargem{src/cryptonote_core/cryptonote_tx_utils.cpp}{cryptonote_tx_utils.cpp}%
  \codigoMargem{construct_miner_tx()}}%
Uma transacção de mineiro criada à altura \(h\) tem obrigatoriamente
\(\texttt{unlock\_time}=h+60\). A sua saída só pode portanto ser referida a
partir dessa altura; o bloqueio de 60 domina a maturação ordinária de 10.


\subsection{Da igualdade de pontos ao valor conservado}
\label{gamma}
\index{compromisso@compromisso!vinculacao no equilibrio@vinculação no equilíbrio}

A equação~\eqref{eq:transaction-commitment-balance} é a mesma
qualquer que seja o número de entradas. De um lado está uma soma de
compromissos cujos montantes vêm de saídas anteriores, do outro a soma
dos compromissos novos e a taxa. Denotemos por \(d\) a diferença das
duas somas inteiras de montantes e por \(u\) a diferença das somas
de máscaras. O significado da comparação é
\[
                            uG+dH=\mathcal O.
\]
As máscaras honestas dão \(u=0\). Se alguém apresentasse aberturas
conhecidas para as quais \(d\ne0\pmod l\), a igualdade daria
\(H=-u d^{-1}G\). É neste passo que intervém a hipótese de vinculação.
A extracção das aberturas nas provas e a dificuldade de encontrar a
relação entre \(G\) e \(H\) permitem interpretar a igualdade como
\(d=0\pmod l\).

As saídas ordinárias admitem montantes em \([0,2^{64})\), e o seu
número por transacção está limitado. Do lado das entradas, a
assinatura liga cada montante a uma saída válida anterior, e o peso
da transacção limita o número de referências que ela pode conter.
Mesmo a cota grosseira de uma entrada por byte, no limite de
\(10^6\) bytes da serialização considerada, deixa ambas as somas
muito aquém de \(l>2^{252}\). Assim, dois totais diferentes dentro
dos limites admitidos têm resíduos diferentes. A conclusão acerca
dos inteiros exige tanto estes limites como as provas de intervalo.
No capítulo~\ref{chapter:dlog-relations}, a mesma equação reaparecerá
na comparação de duas aberturas de uma família de bases.

\subsection{Três entradas para o mesmo pagamento}
\label{sec:tres-entradas-equilibrio}
\index{pseudo-compromisso@pseudo-compromisso!escolha das mascaras@escolha das máscaras}

No princípio Alice tinha as onze unidades numa só saída. Imaginemos
que as recebeu em três ocasiões diferentes, 2 da primeira vez, 4 da
segunda e 5 da terceira. Bob há-de receber os mesmos 7, o troco
continua a ser 3 e a taxa 1. O que muda na conta do nó? Há agora
três entradas, cada qual com o seu anel e a sua imagem de chave.
No entanto, a soma dos montantes de entrada só tem de satisfazer
uma igualdade. Esta diferença entre três autorizações e uma única
conta explica a liberdade que a carteira tem ao escolher as máscaras.

Para acompanhar a álgebra até ao fim, demos valores pequenos às
máscaras do exemplo. Os geradores são \(G,H\) no subgrupo de ordem
\(l\). Os números que seguem servem apenas para expor as operações.
Alice recebeu as três saídas com os compromissos
\[
       C_1=4G+2H,\qquad C_2=9G+4H,\qquad C_3=12G+5H .
\]
Nas saídas novas suponhamos que as máscaras derivadas para Bob e
para o troco valem 6 e 8. Do lado que paga, o nó terá então de fazer
a soma
\[
               (6G+7H)+(8G+3H)+H=14G+11H .
\]
Se somássemos os compromissos recebidos por Alice, obteríamos
\(25G+11H\). Ambas as expressões representam as onze unidades, e
no entanto há uma diferença de \(11G\). As máscaras nas saídas
antigas nasceram em pagamentos independentes, cujos autores
desconheciam o que Alice havia de fazer hoje.

Chamemos \(u_1,u_2,u_3\) às três máscaras dos pseudo-compromissos.
Para comparar as somas basta que satisfaçam, em \(\mathbb F_l\),
\[
                              u_1+u_2+u_3=14 .
\]
Por exemplo, escolhendo \(u_1=3\) e \(u_2=5\), resta
\(u_3=14-3-5=6\). Ficam \(\widetilde C_1=3G+2H\),
\(\widetilde C_2=5G+4H\) e \(\widetilde C_3=6G+5H\).
As três autorizações usam, nas respectivas posições secretas, as
diferenças que se podem ler na última coluna.
\[
\begin{array}{c|c|c|c}
 j & C_j &\widetilde C_j&C_j-\widetilde C_j \\ \hline
 1&4G+2H &3G+2H&G\\
 2&9G+4H &5G+4H&4G\\
 3&12G+5H&6G+5H&6G
\end{array}
\]
Na primeira entrada a testemunha auxiliar de CLSAG vale 1, na
segunda vale 4 e na terceira vale 6. Cada uma acompanha a chave
privada de gasto da sua própria saída. Vejamos o que o nó consegue
refazer desta tabela. Para cada um dos dezasseis membros de um anel
conhece um compromisso histórico e pode subtrair-lhe o
pseudo-compromisso dessa entrada. A assinatura atesta o conhecimento
das duas testemunhas na mesma posição. A coluna que aqui exibe os
três resultados reais esteve à nossa disposição porque acompanhámos
as escolhas da carteira. Ao nó chegam três listas de dezasseis
diferenças, juntamente com a autorização de cada uma.

Para conferir o balanço, o nó já dispõe dos três
pseudo-compromissos sem procurar membro algum dos anéis. A soma é
\[
 (3G+2H)+(5G+4H)+(6G+5H)
      =(3+5+6)G+(2+4+5)H=14G+11H .
\]
É o ponto que obteve do lado das novas saídas e da taxa. As imagens
de chave têm outra função, que se conserva mesmo se duas listas de
dezasseis membros partilharem alguns candidatos. O nó verifica as
três imagens e o seu estado de utilização. Uma coincidência entre
candidatos de dois anéis é possível. Cada entrada tem a sua imagem
e o seu fecho.

Vale a pena contar o que a igualdade das máscaras determina.
Tendo \((3,5,6)\) por solução, todo o triplo
\[
                  (3+s,\ 5+t,\ 6-s-t),\qquad
                  s,t\in\mathbb F_l ,
\]
dá a mesma soma. Reciprocamente, quaisquer máscaras com soma 14
determinam \(s=u_1-3\), \(t=u_2-5\), e a última componente fica
fixada. Eis uma correspondência entre os pares de escalares e as
soluções da equação. A equação linear, considerada em todo o corpo,
tem por isso \(l^2\) soluções. Uma regra concreta de amostragem
pode restringir os valores sorteados. Para conhecer a distribuição
com que a carteira percorre as soluções é preciso conhecer também
a lei desses sorteios.

A liberdade tem uma interpretação pública. Se passarmos de
\((3,5,6)\) para \((5,5,4)\), o primeiro pseudo-compromisso recebe
\(2G\), o último perde \(2G\), e a soma fica intacta. Em contrapartida,
as listas de diferenças dos anéis correspondentes mudam, e cada
assinatura tem de ser construída para o pseudo-compromisso efectivamente
declarado. A igualdade dos totais permite várias escolhas antes da
assinatura. A mensagem e o transcrito fixam as escolhas efectuadas
para aquela transacção.

Com \(m\) entradas, a mesma conta deixa \(m-1\) coordenadas livres.
Fixada a soma \(Y\) das máscaras das saídas, a última coordenada é
\[
                 u_m=Y-\sum_{j=1}^{m-1}u_j .
\]
A uma entrada, a soma vazia vale zero, o que explica por que razão a
carteira de Alice, no primeiro exemplo, tinha uma única máscara possível para
o pseudo-compromisso. A três entradas tem duas escolhas preliminares.
Em ambos os casos a última vem da mesma igualdade. Cada entrada
precisa da sua autorização. Passar de uma para três pede
mais duas assinaturas e mais duas imagens, ainda que os montantes
entregues a Bob e a Alice fiquem exactamente nos mesmos compromissos.
A prova de intervalo abrange as duas saídas novas.


\section{Resumo conceptual}
\label{sec:transaction_summary}

Reunamos os elementos de uma transacção v16. Quem quiser acompanhar
a sua disposição em bytes pode seguir um exemplar do tipo 6 no
apêndice~\ref{appendix:pagamento-tipo-seis}. O apêndice~\ref{appendix:RCTTypeBulletproof2}
mostra como se dispõem os campos numa transacção mais antiga.

\begin{itemize}
    \item \underline{Tipo}: `6' é \texttt{RCTType}\allowbreak\texttt{BulletproofPlus}; `0' é
    \texttt{RCTTypeNull} para a transacção de mineiro.

    \item \underline{Entradas}: para cada entrada \(j\):
    \begin{itemize}
        \item montante visível zero;
        \item dezasseis \emph{offsets} relativos que localizam o anel;
        \item imagem de chave \(I_j\);
        \item pseudo-compromisso \(\widetilde C^a_j\), na parte podável;
        \item CLSAG com dezasseis \(s_{i,j}\), \(c_{1,j}\) e
        \(\bar D_j\), também na parte podável.
    \end{itemize}

    \item \underline{Saídas}: para cada saída \(t\):
    \begin{itemize}
        \item montante visível zero;
        \item alvo etiquetado com o endereço oculto \(P^b_t\) e um
        \emph{view tag} de um byte;
        \item compromisso normal \(C^b_t=y_tG+b_tH\) em \texttt{outPk};
        \item oito bytes \(b_t\) cifrados por XOR em \texttt{ecdhInfo}.
    \end{itemize}

    \item \underline{Base RingCT}: tipo, taxa pública, \texttt{ecdhInfo} e
    \texttt{outPk}. Esta parte não é podável.

    \item \underline{Parte podável}: prova Bulletproofs+, CLSAG e
    pseudo-compromissos.

    \item \underline{\texttt{extra}}: bytes opacos no prefixo, onde as
    carteiras colocam normalmente a chave pública principal de transacção,
    eventuais chaves adicionais e, quando aplicável, um \emph{nonce}. O
    consenso assina os bytes exactos, mas não prova que estes obedecem às
    relações secretas pretendidas pela carteira.
\end{itemize}

Numa transacção com uma entrada e duas saídas, a validade segundo o consenso e
a utilização pelos destinatários dependem das condições seguintes:
\begin{itemize}
 \item para que a transacção seja utilizável, as chaves públicas de transacção
 permitam aos destinatários reconhecer as respectivas saídas;
 \item a CLSAG prove o conhecimento, na mesma posição secreta do anel, da
 chave privada da saída e da diferença de máscaras entre o compromisso
 histórico e o pseudo-compromisso;
 \item a imagem de chave ainda não tenha aparecido num gasto aceite;
 \item cada saída inclua endereço oculto, etiqueta de vista, compromisso e
 montante cifrado;
 \item a Bulletproofs+ prove que os dois montantes pertencem ao intervalo de
 64 bits;
 \item a taxa \(f\) seja pública e se verifique
 \(\widetilde C^a=C^b_0+C^b_1+fH\); e
 \item a mensagem CLSAG vincule o prefixo, a base RingCT e a prova, enquanto a
 transcrição CLSAG vincula também o pseudo-compromisso.
\end{itemize}
Estas condições ocultam a posição real do anel e os montantes transferidos,
sem impedir a verificação da autorização, da ausência de gasto duplo, do
intervalo dos montantes e do equilíbrio.


\subsection{Requisitos de espaço}
\label{subsec:space-and-ver-rcttypefull}
\index{transaccao@transacção!tamanho e peso@tamanho e peso}

Um anel v16 não repete no \emph{blob} os 16 endereços ocultos e os 16
compromissos: os \emph{offsets} apontam para esses dados na lista de blocos. Por
entrada, os elementos fixos directamente ligados ao gasto são:
\begin{align*}
 \underbrace{16\cdot32+32+32}_{\text{CLSAG}}
 +\underbrace{32}_{I_j}
 +\underbrace{32}_{\widetilde C^a_j}
 =640\ \text{bytes},
\end{align*}
antes dos dezasseis \emph{varints} dos \emph{offsets} e dos pequenos
prefixos de serialização. A CLSAG propriamente dita ocupa 576 bytes: 512 nas
respostas, 32 em \(c_1\) e 32 em \(\bar D\).

Para as saídas, cada \(t\) acrescenta pelo menos os 32 bytes do endereço
oculto, um byte de \emph{view tag}, 32 bytes do compromisso e oito bytes do
montante em código, mais as marcas e comprimentos da serialização. As chaves
de transacção vivem em \texttt{extra}. A prova Bulletproofs+ custa
\begin{align*}
 32\bigl(6+2\log_2(64M)\bigr),
 \qquad M=2^{\lceil\log_2p\rceil},
\end{align*}
mais dois comprimentos de vector.

\subsubsection*{Peso, não apenas bytes}

\notaFonte{%
  \fonteHistorica{src/cryptonote_core/tx_pool.cpp}{tx_pool.cpp}%
  \codigoMargem{get_transaction_weight_limit()}%
  \fonteMargem{src/cryptonote_basic/cryptonote_format_utils.cpp}{cryptonote_format_utils.cpp}%
  \codigoMargem{get_transaction_weight()}\par \codigoMargem{get_transaction_weight_clawback()}%
  \fonteMargem{src/cryptonote_core/tx_verification_utils.cpp}{tx_verification_utils.cpp}%
  \codigoMargem{get_transaction_weight_limit()}}%
Para este formato de transacção, o peso é o tamanho serializado mais uma
penalização de ``recuperação'' da economia obtida pela agregação de
Bulletproofs+. Se \(M\leq2\), a penalização é zero e peso e bytes coincidem.
Para \(M>2\), a regra de peso é
\begin{align*}
 w={}&|\mathrm{blob}|+
 \frac{4}{5}\left[
 320M-32\bigl(6+2\log_2(64M)\bigr)
 \right],
\end{align*}
com divisão inteira no último termo. Na versão 16 uma transacção não pode
exceder
\begin{align*}
 \frac{300\,000}{2}-600=149\,400
\end{align*}
unidades de peso. A reserva de 600 deixa espaço para a transacção de mineiro;
a relação com o peso dinâmico do bloco é explicada na secção
\ref{subsec:dynamic-block-weight}.

\subsubsection*{Quanto acrescenta a saída seguinte?}
\index{peso@peso!agregacao das saidas@agregação das saídas}
\index{prova de intervalo@prova de intervalo!dimensao e numero de saidas@dimensão e número de saídas}

As contagens anteriores permitem responder a uma pergunta de carteira.
Partimos de duas saídas. Se for preciso fazer um terceiro pagamento
na mesma transacção, quantos dados novos surgem do lado das provas?
E que diferença faria passar de três saídas para quatro?

Escrevamos \(p\) para o número efectivo de saídas e \(M\) para a
menor potência de dois que contém \(p\). A prova trabalha com
\(64M\) posições e vai reduzindo essa dimensão à metade. Cada
duplicação de \(M\) exige mais uma ronda, e cada ronda acrescenta
dois pontos de 32 bytes. No intervalo escolhido, os dois comprimentos
dos vectores da prova continuam a ocupar um byte cada um. O número
de provas agregadas também se conserva. Aplicando a fórmula aos
quatro corpos de prova, obtêm-se os comprimentos do quadro.
\[
\begin{array}{r|r|r|r|r}
p&M&64M&\text{rondas}&\text{corpo da prova (bytes)}\\ \hline
2 &2&128 &7 & 32(6+14)=640\\
3 &4&256 &8 & 32(6+16)=704\\
4 &4&256 &8 & 704\\
5 &8&512 &9 & 32(6+18)=768
\end{array}
\]
A terceira saída custa 64 bytes adicionais \emph{na prova}. A
quarta ocupa a capacidade da mesma dimensão e acrescenta zero bytes
nesta coluna. Ainda terá de deixar no prefixo a sua própria chave
e etiqueta, e na base RingCT o seu compromisso e o seu montante
cifrado. Também a verificação usa os dados de todas as saídas que
a declaração contém.

Podemos acabar a conta dos bytes com hipóteses bastante precisas.
Conservemos o número de entradas, o tamanho dos seus deslocamentos,
a mesma taxa e o mesmo comprimento de \texttt{extra}. Pensemos em
destinos com endereços principais, para os quais chega a mesma
chave de transacção. Uma nova saída traz um byte de montante visível
zero, um de tipo de alvo, 32 de chave e um de etiqueta. São
\(1+1+32+1=35\) no prefixo. Na base RingCT há ainda oito bytes
de montante cifrado e 32 de compromisso. Para as contagens da
tabela, o comprimento de \texttt{vout} cabe sempre num byte.
Assim, a despesa fixa por saída é \(35+8+32=75\) bytes. A
passagem de duas para três saídas aumenta a serialização em
\(75+64=139\) bytes, e a de três para quatro em 75.
Nesta comparação medimos comprimentos. Mudar os destinatários
obriga naturalmente a calcular de novo os valores das assinaturas
e dos compromissos.

Podemos agora apreciar a diferença entre este custo e o peso.
Para \(M=2\), a parcela de recuperação é zero. Para \(M=4\)
já vale
\[
 \left\lfloor\frac45(4\cdot320-704)\right\rfloor
       =\lfloor460{,}8\rfloor=460 .
\]
Logo, sob as mesmas hipóteses, acrescentar a terceira saída aumenta
o peso em \(139+460=599\), enquanto acrescentar a quarta o aumenta
em 75. O salto repete-se quando a quinta saída requer \(M=8\).
A parcela de recuperação passa então a
\(\lfloor4(2560-768)/5\rfloor=1433\), um acréscimo de 973.
A quinta saída acrescenta 139 bytes à serialização e
\(139+973=1112\) unidades ao peso. A partir do número de saídas,
o leitor consegue assim antecipar em que degraus muda a dimensão
da prova e em quais muda somente o número de objectos comprometidos.

Há dois números, pois, para responder à pergunta sobre a saída
seguinte. Os bytes contam o que precisa de ser transportado e
guardado na forma integral. O peso leva essa contagem à regra
segundo a qual a transacção ocupa espaço no bloco. Uma chave
adicional em \texttt{extra}, ou uma taxa cujo \emph{varint} cresça,
acrescenta os seus próprios bytes ao cálculo. É o campo
\texttt{extra} que precisamos agora de abrir.


\subsection{O campo \texttt{extra}: leitura e significado}
\label{subsec:tx-extra-structure}
\index{extra@\texttt{extra}}
\index{tx extra@\texttt{tx\_extra}|see{extra}}
\index{nonce@nonce!campo auxiliar da transaccao@campo auxiliar da transacção}

O prefixo termina com \texttt{extra}, um vector de bytes que contém
informação auxiliar para as carteiras. A expressão habitual
\texttt{tx\_extra} designa este campo da transacção. O seu comprimento
\(e\) é escrito como \emph{varint}, seguido de exactamente \(e\) bytes.
Fica, portanto, depois de \texttt{vin} e \texttt{vout}. Os
\texttt{key\_offsets}, já lidos dentro de cada entrada de \texttt{vin},
localizam os membros dos anéis. As chaves em \texttt{extra} ajudam o
destinatário a reconhecer e decifrar as novas saídas.

As convenções de leitura distinguem os campos interiores por uma etiqueta
de um byte. A etiqueta determina a forma do conteúdo que se segue:
\begin{center}
\begin{tabular}{lll}
\hline
Etiqueta & Conteúdo & Bytes após a etiqueta\\
\hline
\texttt{01} & Chave principal \(R\) & \(\operatorname{enc}(R)\), 32 bytes\\
\texttt{04} & \(q\) chaves adicionais & \(\operatorname{varint}(q)\), depois \(32q\) bytes\\
\texttt{02} & \emph{Nonce} de \(u\) bytes & \(\operatorname{varint}(u)\), depois \(u\) bytes\\
\texttt{00} & Enchimento & Sequência de zeros\\
\hline
\end{tabular}
\end{center}
As marcas estão escritas em hexadecimal. A chave principal tem comprimento
fixo. Para as chaves adicionais, o inteiro conta pontos; para o
\emph{nonce}, conta bytes. O leitor avança segundo essa regra, de modo
que um byte \texttt{01} no interior de uma chave continua a pertencer
à chave. O enchimento reconhecido ocupa o final do vector, com até 255
bytes zero, incluindo a etiqueta. O \emph{nonce} admite até 255 bytes
de conteúdo.

Considere-se uma transacção com duas saídas que publica a chave principal
\(R\), duas chaves adicionais \(R_0,R_1\) e um identificador cifrado de
oito bytes no \emph{nonce}. A organização do conteúdo é
\begin{align*}
 &\texttt{01}\ \|\ \operatorname{enc}(R)\\
 &\texttt{04}\ \|\ \texttt{02}\ \|\
       \operatorname{enc}(R_0)\ \|\ \operatorname{enc}(R_1)\\
 &\texttt{02}\ \|\ \texttt{09}\ \|\
       \texttt{01}\ \|\ \text{oito bytes cifrados}.
\end{align*}
Na segunda linha, \texttt{02} conta as duas chaves. Na terceira,
\texttt{09} conta nove bytes de conteúdo: a marca interior \texttt{01},
que identifica esta convenção de identificador cifrado, e os oito bytes
seguintes. O total é
\[
 (1+32)+(1+1+64)+(1+1+9)=33+66+11=110.
\]
O vector recebe, por isso, o prefixo exterior \texttt{6E}, pois
\(110=\texttt{6E}_{16}\). Incluindo esse prefixo, ocupa 111 bytes.
Já um \texttt{extra} que contenha apenas a chave principal ocupa 33
bytes de conteúdo e começa exteriormente por \texttt{21}.

A presença das chaves adicionais segue os casos de destinatários e
sub-endereços da secção~\ref{sec:multi_out_transactions}. Quando se
publica uma chave por saída, a posição na lista corresponde ao índice
local dessa saída. A carteira usa a chave pertinente na derivação
partilhada, confere a etiqueta de vista, reconstrói o endereço oculto e
verifica o compromisso do montante decifrado. Obtém assim uma ligação
operacional entre os bytes auxiliares e os objectos matemáticos da
saída. O exemplo de \emph{nonce} descreve uma convenção de leitura;
a autorização do gasto é fornecida pelas assinaturas.

A carteira de referência ordena campos que consegue analisar e recusa criar
\texttt{extra} com mais de 1060 bytes. Um nó recusa normalmente propagar pela
\emph{mempool} uma transacção que exceda esse valor. O limite de 1060 é
política, não uma regra de consenso do bloco: no consenso, \texttt{extra}
continua limitado indirectamente pelo tamanho e pelo peso de toda a
transacção.

O vector e o seu comprimento pertencem à serialização do prefixo, cujos
bytes exactos entram em \(h_0\). A autorização vincula, portanto, esse
conteúdo. A validação de consenso trata-o como uma sequência de bytes;
a interpretação dos campos reconhecidos e a correspondência entre
chaves públicas, segredos partilhados e destinatários pertencem à
carteira. A utilização de uma saída exige ambas as camadas: uma
transacção admitida e dados que o destinatário consiga reconhecer e
abrir. Esta separação explica por que a integridade assinada de
\texttt{extra} e a correcção das suas relações de derivação são
propriedades distintas~\cite{monero-source-v16}.

Uma vez aceites os dados pela carteira e pelo nó, Bob pode já reconhecer
os seus 7, e Alice os 3 de troco. Para contar uma confirmação, Bob terá
ainda de encontrar a transacção num bloco do ramo que a sua rede aceita.
É aí que a ordem dos gastos adquire significado comum.
